% sec:sr-minkowski — The Minkowski Metric, Four-Vectors, and Index Gymnastics
\section{The Minkowski Metric, Four-Vectors, and Index Gymnastics}
\label{sec:sr-minkowski}

The geometry of flat spacetime is entirely encoded in a single object: the
metric.  Every interval, every dot product, every index-raising operation in
the codebase passes through it.  In curved spacetime the metric will vary from
point to point---that's the content of general relativity---but the algebraic
operations are identical.  Learning them here, where the metric is constant, is
like learning matrix multiplication with numbers before doing it with
functions.

\paragraph{Events and coordinates.}

An \emph{event} is a point in spacetime: a place and a time.  We label events
by four coordinates $x^\mu = (x^0, x^1, x^2, x^3) = (t, x, y, z)$, where the
index $\mu$ runs over $0, 1, 2, 3$.
\physnote{The superscript on $x^\mu$ is an \emph{index}, not an exponent.
Context always makes this clear.}
The coordinates themselves aren't physical---different observers may use
different labels for the same event---but the relationships between events
\emph{are} physical, and the metric is what makes them precise.

\paragraph{The Minkowski metric.}

The \emph{Minkowski metric} $\eta_{\mu\nu}$ assigns a notion of ``interval''
to pairs of nearby events.  In Cartesian coordinates, it takes the form
%
\begin{equation}
  \eta_{\mu\nu}
    = \operatorname{diag}(-1,\, +1,\, +1,\, +1) \,,
  \label{eq:minkowski-metric}
\end{equation}
%
so that the spacetime interval between two infinitesimally separated events is
%
\begin{equation}
  \d s^2
    = \eta_{\mu\nu}\,\d x^\mu\,\d x^\nu
    = -\d t^2 + \d x^2 + \d y^2 + \d z^2 \,.
  \label{eq:minkowski-interval}
\end{equation}
%
The crucial feature is the minus sign on $\d t^2$.  Without it, spacetime would
be four-dimensional Euclidean space and there would be no distinction between
time and space, no light cones, and no causality.  With it, we get the
three-way classification of directions that underpins all of relativistic
physics.
\warnnote{The other common convention is $(+{-}{-}{-})$.  We use $(-{+}{+}{+})$
throughout, following \cite{MTW1973}.  Formulae from sources using the opposite
convention will differ by an overall sign in certain places.}

We've used the \emph{Einstein summation convention}: any index appearing once
as a superscript and once as a subscript in a single term is summed over all
four values $0, 1, 2, 3$.  The expression $\eta_{\mu\nu}\,\d x^\mu\,\d
x^\nu$ therefore stands for the double sum
%
\begin{equation*}
  \eta_{\mu\nu}\,\d x^\mu\,\d x^\nu
    \;\equiv\;
    \sum_{\mu=0}^{3}\sum_{\nu=0}^{3}
    \eta_{\mu\nu}\,\d x^\mu\,\d x^\nu \,.
\end{equation*}
%
We will use this convention without further comment throughout the book.

\paragraph{Contravariant and covariant components.}

A \emph{four-vector} $V$ lives in the tangent space and has
\emph{contravariant} components $V^\mu$ (upper indices).  Its metric dual is a
covector with \emph{covariant} components $V_\mu$ (lower indices).  These are
geometrically distinct objects---the metric provides a canonical isomorphism
between them, so we pass freely between the two.  Lowering an index:
%
\begin{equation}
  V_\mu = \eta_{\mu\nu}\, V^\nu \,.
  \label{eq:index-lowering}
\end{equation}
%
Since $\eta_{\mu\nu}$ is diagonal with entries $(-1, +1, +1, +1)$, this
operation has a simple effect on the components: the spatial components are
unchanged, and the time component picks up a minus sign.  Explicitly,
$V_0 = -V^0$ while $V_i = V^i$ for $i = 1, 2, 3$.

Raising works with the \emph{inverse metric} $\eta^{\mu\nu}$, defined by the
condition $\eta^{\mu\alpha}\,\eta_{\alpha\nu} = \delta^\mu{}_\nu$.  For
Minkowski spacetime in Cartesian coordinates, the inverse happens to have the
same component values as the metric itself:
%
\begin{equation}
  \eta^{\mu\nu}
    = \operatorname{diag}(-1,\, +1,\, +1,\, +1) \,.
  \label{eq:minkowski-inverse}
\end{equation}
%
This is a coincidence of flat spacetime in Cartesian coordinates---for the
Schwarzschild or Kerr metrics, the inverse metric will have very different
components from the metric.  But the operation $V^\mu = \eta^{\mu\nu}\,V_\nu$
has the same structure in all cases: contract the metric (or its inverse) with
the vector.

Think of the metric as a translation dictionary between two different
languages.  Contravariant vectors and covariant covectors are different kinds of
geometric object: a velocity $\d x^\mu / \d\tau$ lives in the tangent space,
while a gradient $\partial_\mu \phi$ lives in the cotangent space.  They share
the same representation---an ordered list of four numbers---but they transform
differently and measure different things.  The metric provides an isomorphism
between the two spaces, converting one into the other freely.  This carries
through to curved spacetime, where the dictionary changes from point to point
but the two-space structure persists.
\xrefnote{In curved spacetime (\cref{ch:differential-geometry}),
$\eta_{\mu\nu}$ generalises to $\metric$ and the components become
functions of position.  The index operations look identical.}

\paragraph{The inner product and vector classification.}

The metric defines an inner product (indefinite, unlike the Euclidean case)
between any two four-vectors $A^\mu$ and $B^\mu$:
%
\begin{equation}
  A \cdot B
    = \eta_{\mu\nu}\, A^\mu\, B^\nu
    = -A^0 B^0 + A^1 B^1 + A^2 B^2 + A^3 B^3 \,.
  \label{eq:minkowski-inner-product}
\end{equation}
%
Unlike the Euclidean dot product, this can be negative, zero, or positive.  The
sign of a vector's ``squared norm'' $V^2 \equiv \eta_{\mu\nu}\,V^\mu\,V^\nu$
classifies it into one of three types:
%
\begin{itemize}
  \item \textbf{Timelike} ($V^2 < 0$): the vector points more into the future
    (or past) than into space.  Massive particles travel along timelike
    worldlines.
  \item \textbf{Null} ($V^2 = 0$): the vector lies exactly on the light cone.
    Photons travel along null geodesics---these are the trajectories that the
    ray tracer computes, millions of them per image.
  \item \textbf{Spacelike} ($V^2 > 0$): the vector points more into space
    than into time.  No massive particle or photon can travel along a spacelike
    direction.
\end{itemize}
%
This classification is \emph{invariant}: all observers agree on whether a
vector is timelike, null, or spacelike, regardless of their relative motion.
The metric makes this possible by providing an observer-independent notion of
``norm''.
\physnote{The null condition $V^2 = 0$ doesn't mean the vector is zero---it
means the time and space contributions cancel exactly.  A photon moves through
space precisely as fast as it moves through time.}

\begin{figure}[htbp]
  \centering
  % minkowski-lightcone.tex — Light cone and causal vector classes in 1+1 Minkowski space
% Inputable TikZ fragment for fig:minkowski-lightcone
% Requires: tikz with calc, arrows.meta (loaded by ehbook.cls)
%
\begin{tikzpicture}[
    >={Stealth[length=5pt]},
    font=\footnotesize,
  ]

  \def\axlen{2.8}     % axis half-length
  \def\cone{2.4}      % light-cone extent along each ray

  % ── Shaded future / past regions ──────────────────────────────────
  \fill[EHAccretionGold, fill opacity=0.08]
    (0,0) -- (-\cone,\cone) -- (\cone,\cone) -- cycle;
  \fill[EHJetBlue, fill opacity=0.06]
    (0,0) -- (-\cone,-\cone) -- (\cone,-\cone) -- cycle;

  % Region labels
  \node[font=\scriptsize, text=EHMarginGray] at (0, 1.7) {future};
  \node[font=\scriptsize, text=EHMarginGray] at (0,-1.7) {past};

  % ── Light-cone edges ──────────────────────────────────────────────
  \draw[EHMarginGray, thin] (-\cone,\cone)  -- (0,0) -- (\cone,\cone);
  \draw[EHMarginGray, thin] (-\cone,-\cone) -- (0,0) -- (\cone,-\cone);

  % ── Axes ──────────────────────────────────────────────────────────
  \draw[->, thick, EHMarginGray] (-\axlen,0) -- (\axlen,0)
    node[right, font=\small] {$x$};
  \draw[->, thick, EHMarginGray] (0,-\axlen) -- (0,\axlen)
    node[above, font=\small] {$t$};

  % ── Timelike vector (gold, inside future cone) ────────────────────
  \draw[->, very thick, EHAccretionGold]
    (0,0) -- (0.45,1.4)
    node[left=2pt, text=EHAccretionGold!85!black] {timelike};

  % ── Spacelike vector (blue, outside cone) ─────────────────────────
  \draw[->, very thick, EHJetBlue]
    (0,0) -- (1.9,0.4)
    node[right=2pt, text=EHJetBlue!85!black] {spacelike};

  % ── Null vector (crimson, along cone edge) ────────────────────────
  \draw[->, very thick, EHHorizonCrimson]
    (0,0) -- (1.55,1.55)
    node[right=2pt, text=EHHorizonCrimson!85!black] {null};

\end{tikzpicture}

  \caption[Causal classes of four-vectors in Minkowski space]{Spacetime diagram of 1+1 Minkowski space, showing the three causal
    classes of four-vectors.  Timelike vectors (gold) lie inside the light cone
    and have negative norm $\eta_{\mu\nu} V^\mu V^\nu < 0$; spacelike vectors
    (blue) lie outside and have positive norm; null vectors (crimson) lie
    exactly on the cone with vanishing norm.  The photon trajectories computed
    by the ray tracer are always null---they live on this cone and its
    curved-spacetime generalisation.}
  \label{fig:minkowski-lightcone}
\end{figure}

\paragraph{A worked example.}

Consider a photon travelling in the $x$-direction.  Over a coordinate time
interval $\Delta t$, its displacement four-vector is
%
\begin{equation}
  \Delta x^\mu = (\Delta t,\, \Delta t,\, 0,\, 0) \,,
  \label{eq:photon-displacement}
\end{equation}
%
where we use $c = 1$.  Contract with the metric:
%
\begin{equation}
  \Delta s^2
    = \eta_{\mu\nu}\,\Delta x^\mu\, \Delta x^\nu
    = -(\Delta t)^2 + (\Delta t)^2 + 0 + 0
    = 0 \,.
  \label{eq:photon-interval}
\end{equation}
%
The interval vanishes: the displacement is null, exactly as required for a
photon.  This isn't a derived result---it's the \emph{definition} of what it
means to travel at the speed of light.  The entire ray tracer is built on
propagating this constraint through curved spacetime, where the metric changes
from $\eta_{\mu\nu}$ to $\metric$ and the algebra becomes a system of
differential equations.

For a spacelike example, take two events at the same coordinate time separated
by a distance $d$: $\Delta x^\mu = (0,\, d,\, 0,\, 0)$.  Then $\Delta s^2 =
d^2 > 0$---spacelike, as expected for a purely spatial separation.

Now consider a massive particle moving in the $x$-direction at speed $v < 1$.
Over a coordinate time $\Delta t$, the displacement is $\Delta x^\mu = (\Delta
t,\, v\,\Delta t,\, 0,\, 0)$, giving
%
\begin{equation}
  \Delta s^2
    = -(1 - v^2)\,(\Delta t)^2
    < 0 \,.
  \label{eq:timelike-interval}
\end{equation}
%
The interval is negative: the displacement is timelike.  For any timelike
displacement ($\Delta s^2 < 0$), the quantity
$\Delta\tau = \sqrt{-\Delta s^2}$ is the proper time elapsed on the particle's
own clock---less than the coordinate time $\Delta t$ for any $v > 0$, and
equal to it in the rest frame.

\paragraph{Looking ahead.}

The single minus sign in $\eta_{\mu\nu}$ is the source of all causal
structure: it splits spacetime into timelike, null, and spacelike directions,
gives us light cones, and makes the invariant interval possible.  Without it,
there is no distinction between time and space---and no ray tracing.

Everything we've done in this section generalises directly.  When we move to
curved spacetime in \cref{ch:differential-geometry}, the constant matrix
$\eta_{\mu\nu}$ is replaced by a position-dependent tensor field $\metric$,
the interval becomes $\d s^2 = \metric\,\d x^\mu\,\d x^\nu$, and the index
operations acquire exactly the same form but with $\metric$ and $\inversemetric$
in place of $\eta_{\mu\nu}$ and $\eta^{\mu\nu}$.  The classification of
vectors into timelike, null, and spacelike is unchanged---only the specific
components of the metric are different.  This is why we invest in learning the
index machinery here: it's the same machinery everywhere, and once it's
automatic, the physics of curved spacetime is free to take centre stage.
